\begin{tikzpicture}[x=1cm,y=1cm]
	\tikzset{
		zone/.style={draw=Ink!35, line width=0.55pt},
		zonelabel/.style={font=\bfseries\small, text=Ink},
		asset/.style={fig/block, minimum width=30mm, minimum height=9mm, font=\scriptsize}
	}

	\draw[zone, fill=Good!12] (0,0) circle (5.2cm);
	\draw[zone, dashed] (0,0) circle (3.6cm);
	\draw[zone, fill=Warn!13] (0,0) circle (3.6cm);
	\draw[zone, dashed] (0,0) circle (2.1cm);
	\draw[zone, fill=Bad!12] (0,0) circle (2.1cm);

	\node[asset, fill=white] at (0,0.8) {Платёжный сервер};
	\node[asset, fill=white, minimum width=26mm] at (-1.1,-0.7) {БД с PAN};
	\node[asset, fill=white, minimum width=22mm] at (1.2,-0.7) {Терминалы};

	\node[asset, minimum width=22mm] at (-2.8,2.2) {DNS};
	\node[asset, minimum width=22mm] at (2.7,2.0) {NTP};
	\node[asset, minimum width=27mm] at (-3.0,-2.2) {Мониторинг};
	\node[asset, minimum width=24mm] at (2.8,-2.0) {Файрвол};

	\node[asset, minimum width=24mm] at (-4.8,4.0) {CRM};
	\node[asset, minimum width=29mm] at (4.8,4.0) {Маркетинг};
	\node[asset, minimum width=27mm] at (0,-5.0) {HR-системы};

	\node[zonelabel, text=Bad!80!black] at (0,1.85) {CDE};
	\node[zonelabel, text=Warn!70!black] at (0,3.25) {connected-to / security-impacting};
	\node[zonelabel, text=Good!55!black] at (0,4.95) {вне периметра};

	\node[font=\scriptsize, text=Muted, align=center] at (6.9,0.2) {Межсетевой экран /\\сегментация};
	\draw[fig/dashed] (5.8,0.9) -- (3.6,0.9);
	\draw[fig/dashed] (5.8,-0.9) -- (2.1,-0.9);

	\node[fig/note, text width=12.4cm] at (0,-6.9) {Без сегментации соседние инфраструктурные сервисы и бизнес-системы тоже попадают в периметр PCI DSS, поэтому границы CDE всегда нужно проектировать как архитектурное решение, а не как список серверов.};
\end{tikzpicture}
